\appendix

\section{Subspace Reconstruction Identity}
\label{app:subspace_identity}

Let $\mathbf A\in\mathbb R^{C\times r}$ have orthonormal columns and let
$\widehat{\mathbf x}=\mathbf y-\mathbf A\mathbf d$ for any coefficient sequence
$\mathbf d\in\mathbb R^{r\times L}$.  Defining
$\mathbf Q=\mathbf I-\mathbf A\mathbf A^\top$,
\[
  \mathbf Q\widehat{\mathbf x}
  =\mathbf Q\mathbf y-\mathbf Q\mathbf A\mathbf d
  =\mathbf Q\mathbf y,
\]
because $\mathbf Q\mathbf A=\mathbf A-\mathbf A(\mathbf A^\top\mathbf A)=0$.
The implementation evaluates the relative residual of this identity on valid
time samples and requires it to be at most $10^{-5}$.  The identity says only
that the observation is preserved outside the \emph{estimated} artifact basis.

\section{Bounded Coefficients}
For a positive training-only scale $\tau_j$ and physical coefficient $a_j$, the
target $u_j=\tanh(a_j/\tau_j)$ lies in $(-1,1)$.  The sampler applies the same
bound to every predicted coefficient before reconstruction.  Therefore the
per-coordinate physical correction satisfies
$|\tau_j u_j|\leq\tau_j$.  This bound is independent of the evaluation target
and prevents a reverse trajectory from producing an unbounded EEG correction.

\section{Reporting Scope}
Klados evaluation units are source records, not participants.  SGEYESUB units
are participant stems, but all SGE studies were available during development.
The intervals in the main paper are paired-bootstrap development intervals and
must not be interpreted as independent confirmation.  The attempted independent
EEGEyeNet evaluation was blocked before test outcomes were opened.
